%% Elevark fra mattelab.no – laget automatisk av bygg.py. Rediger gjerne til eget bruk.
%% Kompiler med lualatex.
\documentclass[11pt,a4paper]{article}
\usepackage[norsk]{babel}
\usepackage[a4paper,margin=2.2cm,top=2cm,bottom=2.4cm]{geometry}
\usepackage{fontspec}
\usepackage{unicode-math}
\setmainfont{TeX Gyre Pagella}
\setsansfont{TeX Gyre Heros}[Scale=MatchLowercase]
\setmathfont{TeX Gyre Pagella Math}
\usepackage{xcolor}
\definecolor{blekk}{HTML}{16213A}
\definecolor{dempet}{HTML}{566077}
\definecolor{rute}{HTML}{CFC9B8}
\definecolor{m1}{HTML}{2F4AB0}
\definecolor{m2}{HTML}{99266D}
\usepackage[most]{tcolorbox}
\usepackage{qrcode}
\usepackage{needspace}
\usepackage{fancyhdr}
\usepackage[hidelinks]{hyperref}
\pdfvariable suppressoptionalinfo \numexpr 1+2+4+8+16+32+64+128+256+512\relax
\setlength{\parindent}{0pt}
\setlength{\parskip}{0.4em}
\pagestyle{fancy}
\fancyhf{}
\renewcommand{\headrulewidth}{0pt}
\fancyfoot[L]{\sffamily\footnotesize\color{dempet}mattelab.no · Andregradsulikheter · elevark uten fasit}
\fancyfoot[R]{\sffamily\footnotesize\color{dempet}Side \thepage}

\newcommand{\fase}[2]{\tcbox[on line, colframe=#2, colback=white, coltext=#2, boxrule=0.6pt, arc=2pt,
  boxsep=0pt, left=3pt, right=3pt, top=1.5pt, bottom=1.5pt, fontupper=\sffamily\bfseries\footnotesize]{#1}}
\newcommand{\steg}[4]{%
  \par\noindent{\color{rute}\rule{\linewidth}{0.4pt}}\par\nopagebreak
  \noindent\begin{minipage}[t]{2.1cm}\vspace{0pt}\fase{#1}{#2}\end{minipage}%
  \begin{minipage}[t]{\dimexpr\linewidth-2.1cm\relax}\vspace{0pt}#3\end{minipage}\par
  \vspace{#4}}
\newcommand{\del}[4]{%
  \par\noindent{\color{rute}\rule{\linewidth}{0.4pt}}\par\nopagebreak
  \noindent\begin{minipage}[t]{1.1cm}\vspace{0pt}\fase{#1}{#2}\end{minipage}%
  \begin{minipage}[t]{\dimexpr\linewidth-1.1cm\relax}\vspace{0pt}#3\end{minipage}\par
  \vspace{#4}}
\newcommand{\oppgave}[2]{\par\Needspace{12\baselineskip}\vspace{0.6em}%
  {\fontsize{14}{17}\selectfont\bfseries{\color{m1}#1}\enspace #2}\par\vspace{0.2em}}

\begin{document}
\begin{minipage}[t]{0.72\linewidth}\vspace{0pt}
{\sffamily\small\color{dempet}Elevark · 1T · mattelab.no}\par\vspace{0.3em}
{\fontsize{24}{28}\selectfont\bfseries Andregradsulikheter}\par\vspace{1em}
{\sffamily\small Navn: \rule{4.6cm}{0.4pt}\quad Klasse: \rule{1.5cm}{0.4pt}\quad Dato: \rule{1.9cm}{0.4pt}}
\end{minipage}\hfill
\begin{minipage}[t]{0.22\linewidth}\vspace{0pt}\raggedleft
\qrcode[height=2.3cm]{https://mattelab.no/utforskere/andregradsulikheter.html}\par\vspace{0.2em}
{\sffamily\scriptsize\color{dempet}Utforskeren}
\end{minipage}

\vspace{1em}
{\sffamily\small\color{dempet}Utforskeren ligger på \url{https://mattelab.no/utforskere/andregradsulikheter.html}}\par
Hver oppgave har tre steg. \textbf{Gjett} først, uten å bruke utforskeren. \textbf{Test} så i utforskeren. \textbf{Forklar} til slutt med egne ord hvorfor det blir slik.

\oppgave{1}{Over x-aksen}
\steg{Gjett}{m1}{Nullpunktene til $f(x) = x^2 - 2x - 3$ er $x = -1$ og $x = 3$. For hvilke $x$ er $f(x) > 0$? Skisser grafen.}{2.2cm}
\steg{Test}{m2}{Sett opp ulikheten. Sammenlign skissen og svaret med grafen og fortegnsskjemaet.}{2.6cm}
\steg{Forklar}{blekk}{Lag fortegnsskjemaet for $(x + 1)(x - 3)$ selv. Forklar hvorfor produktet er negativt mellom nullpunktene.}{3.4cm}

\oppgave{2}{Er endepunktene med?}
\steg{Gjett}{m1}{Hva er løsningen av $x^2 - 2x - 3 \leq 0$? Hvordan skiller den seg fra løsningen av $x^2 - 2x - 3 < 0$?}{2.2cm}
\steg{Test}{m2}{Velg ≤, og bytt så til <. Se på prikkene i nullpunktene og på intervallet.}{2.6cm}
\steg{Forklar}{blekk}{Hvorfor er $x = 3$ med i den ene løsningen, men ikke i den andre? Skriv begge løsningene som intervall.}{3.4cm}

\oppgave{3}{Negativ $a$}
\steg{Gjett}{m1}{Løs $-x^2 + x + 6 \geq 0$. Tror du svaret ligger mellom nullpunktene eller utenfor?}{2.2cm}
\steg{Test}{m2}{Sett opp ulikheten. Se på raden for $a$ i fortegnsskjemaet. Sett så opp $x^2 - x - 6 \leq 0$ og sammenlign.}{2.6cm}
\steg{Forklar}{blekk}{Hvorfor har de to ulikhetene samme løsning? Hva skjer med ulikhetstegnet når du ganger med $-1$?}{3.4cm}

\oppgave{4}{Alle, ingen eller bare ett tall}
\steg{Gjett}{m1}{Hva er løsningen av $x^2 + 2x + 5 > 0$? Og av $x^2 + 2x + 5 < 0$? Hva med $(x - 2)^2 \leq 0$?}{2.2cm}
\steg{Test}{m2}{Sett opp $x^2 + 2x + 5 > 0$ og bytt til <. Sett så opp $x^2 - 4x + 4 \leq 0$.}{2.6cm}
\steg{Forklar}{blekk}{Bruk diskriminanten og formen på grafen til å forklare svarene.}{3.4cm}

\oppgave{5}{Ballen}
\steg{Gjett}{m1}{En ball blir kastet rett opp. Høyden etter $t$ sekunder er $h(t) = -5t^2 + 20t + 1$ meter. Hvor lenge er ballen høyere enn $16$ meter?}{2.2cm}
\steg{Test}{m2}{Skriv ulikheten $h(t) > 16$ med $0$ på høyre side. Sett den opp i utforskeren og les av løsningen.}{2.6cm}
\steg{Forklar}{blekk}{Løs ulikheten ved regning. Hvorfor kan du dele med $-5$, og hva skjer da med ulikhetstegnet?}{3.4cm}

\end{document}
